\SourceLeaf{395}
\[
\begin{gathered}\mathfrak{P}t=\frac1{\surd t}\mathfrak{P}\frac1t,\\[5pt]
\mathfrak{Q}t=\frac1{\surd t}\mathfrak{R}\frac1t,\\[5pt]
\mathfrak{R}t=\frac1{\surd t}\mathfrak{Q}\frac1t\end{gathered}
\]
genügen, worin die Quadratwurzeln mit solchen Zeichen zu nehmen sind, dass
der reelle Theil positiv ist.

Aus dieser und der anderen von ihm aufgezeichneten Eigenschaft derselben
Functionen, dass nämlich
\[
\begin{gathered}\mathfrak{P}t=\mathfrak{Q}(t+i),\\[5pt]
\mathfrak{Q}t=\mathfrak{P}(t+i),\\[5pt]
\mathfrak{R}t=\surd i.\mathfrak{R}(t+i)\end{gathered}
\]
ist, scheint \textsc{Gauss} den folgenden Satz abgeleitet zu haben:]

\begin{itshape}
Es seien $\alpha,\allowbreak \beta,\allowbreak \gamma,\allowbreak \delta$ ganze reelle Zahlen, $\alpha\delta-\beta\gamma=1$,
$\dfrac{\alpha t-\beta i}{\delta+\gamma ti}=t'$.

Wir unterscheiden 6 Fälle, jenachdem nach dem Modulus 2
\[
\begin{array}{r@{\;}c@{\quad}cccccc}
\alpha&\equiv&1&1&1&0&1&0\\
\beta&\equiv&0&1&0&1&1&1\\
\gamma&\equiv&0&0&1&1&1&1\\
\delta&\equiv&1&1&1&1&0&0
\end{array}
\]
Es ist dann
\[
\begin{array}{r@{\;}c@{\quad}cccccc}
h\mathfrak{P}t'&=&\mathfrak{P}t&\mathfrak{Q}t&\mathfrak{R}t&\mathfrak{R}t&\mathfrak{Q}t&\mathfrak{P}t\\
h\mathfrak{Q}t'&=&\mathfrak{Q}t&\mathfrak{P}t&\mathfrak{Q}t&\mathfrak{P}t&\mathfrak{R}t&\mathfrak{R}t\\
h\mathfrak{R}t'&=&\mathfrak{R}t&\mathfrak{R}t&\mathfrak{P}t&\mathfrak{Q}t&\mathfrak{P}t&\mathfrak{Q}t
\end{array}
\]
\[
h=\surd i^\lambda(\delta+\gamma ti)
\]
\end{itshape}

[worin $\lambda$ für die Factoren der drei Functionen $\mathfrak{P},\allowbreak \mathfrak{Q},\allowbreak \mathfrak{R}$ im Allgemeinen verschiedene
Werthe hat.]

\begin{itshape}
Ist hier $t=\dfrac{\surd d+bi}{a}$, $t'=\dfrac{\surd d+b'i}{a'}$,
$-d=bb-ac=b'b'-a'c'$, so geht die Form $(a,\allowbreak b,\allowbreak c)$ in $(a',\allowbreak b',\allowbreak c')$ über durch die Transformation
$\left(\begin{smallmatrix}\delta,\allowbreak &-\beta\\-\gamma,\allowbreak &\alpha\end{smallmatrix}\right)$.

Zusammenhang zwischen den Formen des negativen Determinanten $-p$ und
den summatorischen Functionen.

Sind nämlich die Formen $(a,\allowbreak b,\allowbreak c)$ $(A,\allowbreak B,\allowbreak C)$ aequivalent, so ist die Function $f$
in Betracht zu ziehen wo $ft\equiv fu$ so wol wenn $\dfrac{t-u}{i}$ ganze Zahl als wenn
$t=\dfrac1u$.
\par
Jeder Classe entspricht dann ein bestimmter Werth von $f\dfrac{\surd p+bi}{a}$.
\end{itshape}

\SourceLeaf{396}
\AGMArticle{18}

[Mit dem Algorithmus des arithmetisch-geometrischen Mittels hat \textsc{Gauss}
einen andern in Verbindung gebracht, welcher ebenfalls wie jener von zwei gegebenen
Grössen, die hier $\alpha$ und $\beta$ bezeichnet werden sollen, ausgeht und auf
eine solche Form zurückgeführt werden kann, dass viele Analogien mit jenem
sich zeigen, wenn nämlich
\[
\begin{aligned}
4a'a'&=(a+b)^2,\quad b'b'=ab,\quad\text{u.\,s.\,f.}\\
4\alpha'\alpha'&=(\alpha+\beta)^2,\quad\beta'\beta'=\alpha\beta,\quad\text{u.\,s.\,f.}
\end{aligned}
\]
gesetzt wird. Ebenso wie die bisherigen Untersuchungen durch Einführung der
die Gleichung $aa=bb+cc$ erfüllenden Grösse $c$ bedeutend übersichtlicher
wurden, wird hier eine entsprechende Vereinfachung der Formeln erreicht, wenn
man $\gamma$ durch die Gleichung $\alpha\alpha=\beta\beta+\gamma\gamma$ und $\delta$ durch
\[
\begin{aligned}bc(b\gamma-c\beta)&=ca(a\gamma-c\alpha)\\&=ab(b\alpha-a\beta)\\&=abc\delta\end{aligned}
\]
so wie $\gamma^{\mathrm n}$, $\delta^{\mathrm n}$ durch dieselben Gleichungen, nachdem allen Zeichen der Index $n$
gegeben ist, einführt. Unter den zwischen diesen Grössen bestehenden Relationen
finden die folgenden bei der Untersuchung dieses Algorithmus vielfache Anwendung:
\[
\begin{gathered}\frac{\alpha+\beta}{a+b}=\frac{\gamma-\delta}{c},\\[5pt]
\frac{\alpha-\beta}{a-b}=\frac{\gamma+\delta}{c},\\[5pt]
\frac{\alpha+\gamma}{a+c}=\frac{\beta+\delta}{b},\\[5pt]
\frac{\alpha-\gamma}{a-c}=\frac{\beta-\delta}{c},\end{gathered}
\]
\[
\begin{aligned}
\alpha'&=\frac1{a'}\left(\frac{\alpha+\beta}{2}\right)^2
=\frac1{c'}\left(\frac{\gamma-\delta}{2}\right)^2,\quad
\beta'=\frac1{b'}\alpha\beta\\
\gamma'&=\frac1{c'}\left(\frac{\alpha-\beta}{2}\right)^2
=\frac1{a'}\left(\frac{\gamma+\delta}{2}\right)^2,\quad
\delta'=\frac1{b'}\gamma\delta
\end{aligned}
\]
\[
\begin{aligned}\alpha\alpha+\delta\delta&=\beta\beta+\gamma\gamma
\\&=\frac ab(\alpha\beta+\gamma\delta)\\&=\frac ac(\alpha\gamma-\beta\delta)
\\&=a(\alpha'+\gamma')\end{aligned}
\]
\[
\begin{aligned}\alpha\alpha-\gamma\gamma&=\beta\beta-\delta\delta
\\&=\frac bc(\beta\gamma-\alpha\delta)\\&=\frac ba(\alpha\beta-\gamma\delta)
\\&=b(\alpha'-\gamma')\end{aligned}
\]
\[
\begin{aligned}\alpha\alpha-\beta\beta&=\gamma\gamma-\delta\delta
\\&=\frac ca(\alpha\gamma+\beta\delta)\\&=\frac cb(\beta\gamma+\alpha\delta)
\\&=2c\surd{\alpha'\gamma'}\end{aligned}
\]
Die Grenzwerthe der Glieder in den sieben Reihen von Grössen $a,\allowbreak b,\allowbreak c,\allowbreak \alpha,\allowbreak \beta,\allowbreak \gamma,\allowbreak \delta$
lassen sich auf diese vier zurückführen:

\SourceLeaf{397}
\[
k=\mathrm M(a,b)=\lim a^{\mathrm n}=\lim b^{\mathrm n}
\]
\[
\begin{aligned}
\surd y&=e^{-\frac\pi2\frac{\mathrm M(a,b)}{\mathrm M(a,c)}}
=\lim{}^{2^n}\!\surd{\frac{c^{\mathrm n}}{4a^{\mathrm n}}}\\
&=\frac{c}{4a}\cdot\frac a{a'}\cdot\surd{\frac{a'}{a''}}
\cdot{}^{2^2}\!\surd{\frac{a''}{a'''}}\\&\quad\cdot{}^{2^3}\!\surd{\frac{a'''}{a''''}}
.\;.{}^{2^{n-1}}\!\surd{\frac{a^{\mathrm n-1}}{a^{\mathrm n}}}.\;.
\end{aligned}
\]
\[
\begin{aligned}
\frac xk&=\lim{}^{2^n}\!\surd{\frac{\alpha^{\mathrm n}}{k}}
=\lim{}^{2^n}\!\surd{\frac{\beta^{\mathrm n}}{k}}\\
&=\frac\beta b\cdot\surd{\frac ba\frac\alpha\beta}
\cdot{}^{2^2}\!\surd{\frac{b'}{a'}\frac{\alpha'}{\beta'}}
\\&\quad\cdot{}^{2^3}\!\surd{\frac{b''}{a''}\frac{\alpha''}{\beta''}}
.\;.\;.{}^{2^n}\!\surd{\frac{b^{\mathrm n-1}}{a^{\mathrm n-1}}\frac{\alpha^{\mathrm n-1}}{\beta^{\mathrm n-1}}}.\;.
\end{aligned}
\]
\[
\frac xk\surd y.\eta^{\pm1}
=\lim{}^{2^{n-1}}\!\surd{\left[\surd{\frac{\gamma^{\mathrm n}}{4k}}\pm\surd{\frac{\delta^{\mathrm n}}{4k}}\right]}
\]
Wenn $\alpha$ und $\beta$ und alle Grössen $a^{\mathrm n}$, $b^{\mathrm n}$, positiv sind, so nehmen
$\dfrac{a^{\mathrm n}}{b^{\mathrm n}}$ und $\dfrac{a^{\mathrm n}a^{\mathrm n}}{b^{\mathrm n}\beta^{\mathrm n}}$ von den Werthen
$\dfrac{a'}{b'}$ und $\dfrac{a'a'}{b'\beta'}$ beständig bis zur Einheit ab;
es ergeben also die vorstehenden Ausdrücke einen bestimmten Werth für $x$.
Das Gleiche folgt für $\eta$, wenn $\gamma$, $\delta$ und alle $c^{\mathrm n}$ und $b^{\mathrm n}$ positiv sind, aus
\[
\begin{aligned}
\frac xk y^{\frac12}\eta&=\lim{}^{2^n}\!\surd{\frac{\gamma^{\mathrm n}}{k}}
=\lim{}^{2^n}\!\surd{\frac{\delta^{\mathrm n}}{k}}\\
&=\frac\delta b\cdot\surd{\frac ba\frac\gamma\delta}
\cdot{}^{2^2}\!\surd{\frac{b'}{a'}\frac{\gamma'}{\delta'}}
\\&\quad\cdot{}^{2^3}\!\surd{\frac{b''}{a''}\frac{\gamma''}{\delta''}}
.\;.\;.{}^{2^n}\!\surd{\frac{b^{\mathrm n-1}}{a^{\mathrm n-1}}\frac{\gamma^{\mathrm n-1}}{\delta^{\mathrm n-1}}}.\;.\;.
\end{aligned}
\]
wenn aber $\delta$ negativ ist aus der von \textsc{Gauss} angewandten Substitution
\[
\begin{gathered}
\surd{\frac{-\delta}{\gamma}}
=\mathrm{tang}\,U\cdot\surd{\frac ab}
=\mathrm{tang}\,V\cdot\surd{\frac ba}
\\=\surd{(\mathrm{tang}\,U\cdot\mathrm{tang}\,V)}
=\mathrm{tang}\,U'\\
U+V=2V'
\end{gathered}
\]
\[
\begin{gathered}
\surd{\frac{-\delta'}{\gamma'}}
=\mathrm{tang}\,2U'\cdot\surd{\frac{a'}{b'}}
=\mathrm{tang}\,2V'\cdot\surd{\frac{b'}{a'}}
\\=\surd{(\mathrm{tang}\,2U'\cdot\mathrm{tang}\,2V')}
=\mathrm{tang}\,2U''\\
U'+V'=2V''
\end{gathered}
\]
\[
\begin{gathered}
\surd{\frac{-\delta''}{\gamma''}}
=\mathrm{tang}\,4U''\cdot\surd{\frac{a''}{b''}}
=\mathrm{tang}\,4V''\cdot\surd{\frac{b''}{a''}}
\\=\surd{(\mathrm{tang}\,4U''\cdot\mathrm{tang}\,4V'')}
=\mathrm{tang}\,4U'''\\
U''+V''=2V'''
\end{gathered}
\]
\[
\begin{gathered}
\surd{\frac{-\delta^{\mathrm n}}{\gamma^{\mathrm n}}}
=\mathrm{tang}\,2^nU^{\mathrm n}\cdot\surd{\frac{a^{\mathrm n}}{b^{\mathrm n}}}
=\mathrm{tang}\,2^nV^{\mathrm n}\cdot\surd{\frac{b^{\mathrm n}}{a^{\mathrm n}}}
\\=\surd{(\mathrm{tang}\,2^nU^{\mathrm n}\cdot\mathrm{tang}\,2^nV^{\mathrm n})}
=\mathrm{tang}\,2^nU^{\mathrm n+1}\\
U^{\mathrm n}+V^{\mathrm n}=2V^{\mathrm n+1}
\end{gathered}
\]
weil dann
\[
\begin{aligned}
\sin U^2&=\frac bc\frac{-\delta}{\alpha},\\[4pt]&\quad
\cos U^2=\frac ac\frac\gamma\alpha,\\[4pt]&\quad
aa\sin U^2+bb\cos U^2=ab\frac\beta\alpha\\[4pt]
\sin V^2&=\frac ac\frac{-\delta}{\beta},\\[4pt]&\quad
\cos V^2=\frac bc\frac\gamma\beta,\\[4pt]&\quad
aa\cos V^2+bb\sin V^2=ab\frac\alpha\beta
\end{aligned}
\]
ist, und diese Gleichungen auch gelten, wenn $a^{\mathrm n},\allowbreak b^{\mathrm n},\allowbreak c^{\mathrm n},\allowbreak \alpha^{\mathrm n},\allowbreak \beta^{\mathrm n},\allowbreak \gamma^{\mathrm n},\allowbreak \delta^{\mathrm n},\allowbreak 2^nU^{\mathrm n},\allowbreak 2^nV^{\mathrm n}$
statt $a,\allowbreak b,\allowbreak c,\allowbreak \alpha,\allowbreak \beta,\allowbreak \gamma,\allowbreak \delta,\allowbreak U,\allowbreak V$ gesetzt werden, so dass also
\[
\eta^{\pm\frac12}
=\lim e^{\pm iU^{\mathrm n}}=\lim e^{\pm iV^{\mathrm n}}=e^{\pm iu}
\]
sich ergibt.

\SourceLeaf{398}
Für $\delta=0$ verschwinden alle $\delta^{\mathrm n}$ und es wird
\[
\begin{gathered}\frac\alpha a=\frac\beta b=\frac\gamma c=.\;.
\\={}^{2^n}\!\surd{\frac{\alpha^{\mathrm n}}{a^{\mathrm n}}}
={}^{2^n}\!\surd{\frac{\beta^{\mathrm n}}{b^{\mathrm n}}}
={}^{2^n}\!\surd{\frac{\gamma^{\mathrm n}}{c^{\mathrm n}}}=\frac xk,\\[5pt]
{}^{2^n}\!\surd{\frac{\gamma^{\mathrm n}}{\alpha^{\mathrm n}}\frac{a^{\mathrm n}}{c^{\mathrm n}}}=\eta=1,\\[5pt] u=0\end{gathered}
\]
Vergleicht man die Grenzwerthe $k,\allowbreak y,\allowbreak x,\allowbreak \eta,\allowbreak u$, zu denen man gelangt, wenn
man bei der Bildung des combinirten Algorithmus von den Grössen $a,\allowbreak b,\allowbreak \alpha,\allowbreak \beta$
ausgegangen ist, mit den Grenzwerthen $k',\allowbreak y',\allowbreak x',\allowbreak \eta',\allowbreak u'$, zu denen man gelangt,
wenn man bei solchem Algorithmus von den bestimmten zuvor erhaltenen Grössen
$a',\allowbreak b',\allowbreak \alpha',\allowbreak \beta'$ als Anfangsglieder ausginge, so ersieht man unmittelbar, dass
\[
\begin{gathered}k'=k,\\[5pt] y'=yy,\\[5pt]\frac{x'}{k'}=\frac{xx}{kk},\\[5pt]
\eta'=\eta\eta,\\[5pt] u'=2u\end{gathered}
\]
sein muss und dass durch die nach diesem Gesetze gebildeten Gleichungen die
den $a^{\mathrm n},\allowbreak b^{\mathrm n},\allowbreak \alpha^{\mathrm n},\allowbreak \beta^{\mathrm n}$ entsprechenden Grenzwerthe $k^{\mathrm n},\allowbreak y^{\mathrm n},\allowbreak x^{\mathrm n},\allowbreak \eta^{\mathrm n},\allowbreak u^{\mathrm n}$ sich ergeben.

Aus den Gleichungen von der Form:
\[
\begin{gathered}\frac{c'}{a'}=\frac{a-b}{a+b},\\[5pt]
\frac{c'}{a'}\frac{\gamma'}{\alpha'}=\left(\frac{\alpha-\beta}{\alpha+\beta}\right)^2,\\[5pt]
\frac{c'}{a'}\frac{\alpha'}{\gamma'}=\left(\frac{\gamma-\delta}{\gamma+\delta}\right)^2\end{gathered}
\]
folgt, dass, wenn alle Glieder des Algorithmus positiv werden, $y$, $y\eta$ und
$\dfrac1\eta$ kleiner als die Einheit sind.

Die von \textsc{Gauss} angegebene Methode zur Bestimmung des Grenzwerthes für
$U^{\mathrm n}$ und $V^{\mathrm n}$ und ebenso die Bestimmung von $\dfrac xk$ mit Zuhülfenahme jener Winkel
führt bei Rechnungen mit Zahlen sehr rasch zum Ziele. Weniger bequem
für Zahlenrechnungen sind die obigen Formeln zur Bestimmung eines reellen $\eta$
und des zugehörigen $\dfrac xk$. Dieser Umstand wird die Veranlassung gewesen sein,
wesshalb \textsc{Gauss} den Algorithmus:]
\[
\begin{aligned}
A'&=\frac{A+B}{2},\quad B'=\frac{2ABa'}{b'(A+B)}\\
A''&=\frac{A'+B'}{2},\quad B''=\frac{2A'B'a''}{b'(A'+B')}
\end{aligned}
\]
u.\,s.\,f. [mit dem Grenzwerthe]
\[
\begin{aligned}\frac Hk&=\frac Bb\surd{\frac{bA}{aB}}
\cdot{}^{4}\!\surd{\frac{b'A'}{a'B'}}
\\&\quad\cdot{}^{8}\!\surd{\frac{b''A''}{a''B''}}
\cdot{}^{16}\!\surd{\frac{b'''A'''}{a'''B'''}}.\;.\end{aligned}
\]
[aufgestellt hat, welcher sich auf den obigen zurückführen lässt, wenn man
$A=\alpha$, $B=\beta$ setzt, weil dann
